\documentclass[12pt]{article}
\usepackage[margin=1in]{geometry}
\title{Research with References}
\author{Academic Author}
\begin{document}
\maketitle

\section{Introduction}
Machine learning has transformed many fields. Convolutional neural networks
achieved breakthrough results in image recognition [1]. Transformer architectures
revolutionized natural language processing [2], leading to large language models
that can perform complex reasoning tasks [3].

\section{Related Work}
The challenge of PDF parsing has been studied extensively. Early work focused on
rule-based approaches [4], while more recent methods employ deep learning for
document layout analysis [5]. Our approach combines traditional heuristics with
adaptive thresholds, avoiding the need for training data.

\section{Discussion}
Despite significant progress, PDF-to-text conversion remains challenging due to
the format's low-level representation. As noted by Smith et al. [6], the gap
between visual presentation and logical structure is the fundamental obstacle.

\begin{thebibliography}{9}
\bibitem{lecun} LeCun, Y., Bengio, Y., \& Hinton, G. (2015). Deep learning. \textit{Nature}, 521(7553), 436--444.
\bibitem{vaswani} Vaswani, A., et al. (2017). Attention is all you need. \textit{NeurIPS}, 5998--6008.
\bibitem{brown} Brown, T., et al. (2020). Language models are few-shot learners. \textit{NeurIPS}, 1877--1901.
\bibitem{mao} Mao, S., Rosenfeld, A., \& Kanungo, T. (2003). Document structure analysis algorithms. \textit{Pattern Recognition}, 36(11), 2225--2243.
\bibitem{zhong} Zhong, X., Tang, J., \& Yepes, A. J. (2019). PubLayNet: largest dataset for document layout analysis. \textit{ICDAR}, 1015--1022.
\bibitem{smith} Smith, R. (2007). An overview of the Tesseract OCR engine. \textit{ICDAR}, 629--633.
\end{thebibliography}
\end{document}